\section{截至2026年7月15日的半年报预告增量}

本节是对7月11日基线的增量更新，不重写原有364家A股预告母池和117只估值总账。通过东方财富官方公告接口及公司公告附件，新增归档9只原此前未命中的主题池公司；每只均完成代码、公司名、标题、金额字段和PDF文本校验。由于AkShare的同类端点本轮返回异常结构，采集切换至官方公告接口，该降级路径已写入结构化数据包。

\noindent
\scriptsize
\begin{tabularx}{\textwidth}{L{1.00cm}L{1.25cm}L{0.95cm}R{0.75cm}R{0.75cm}R{0.70cm}R{0.75cm}X}
\toprule
\textbf{代码} & \textbf{公司} & \textbf{公告日} & \textbf{归母中值} & \textbf{扣非中值} & \textbf{H1 EPS} & \textbf{隐含Q2} & \textbf{处置} \\
\midrule
300223 & 北京君正 & 2026-07-13 & 11.8 & 11.5 & 2.44 & 8.6 & 验证观察 \\
600118 & 中国卫星 & 2026-07-13 & 0.3 & 0.3 & 0.02 & 0.8 & 验证观察 \\
002185 & 华天科技 & 2026-07-14 & 8.0 & 2.4 & 0.24 & 7.1 & 一次性排除 \\
002517 & 恺英网络 & 2026-07-14 & 14.3 & 11.7 & 0.67 & 6.5 & 研报待验证 \\
001309 & 德明利 & 2026-07-15 & 61.0 & 60.5 & 27.13 & 27.5 & 兑现待回撤 \\
002156 & 通富微电 & 2026-07-15 & 17.0 & 7.5 & 1.12 & 13.7 & 价格先行 \\
600584 & 长电科技 & 2026-07-15 & 8.6 & 8.2 & 0.44 & 5.7 & 价格先行 \\
600760 & 中航沈飞 & 2026-07-15 & 4.7 & 4.0 & 0.15 & 3.1 & 预减观察 \\
600893 & 航发动力 & 2026-07-15 & 1.5 & 2.0 & 0.04 & 1.4 & 验证观察 \\
\bottomrule
\end{tabularx}
\normalsize

\textit{注：}归母和扣非净利润单位为人民币亿元；隐含Q2为H1归母净利中值减2026Q1归母净利，不是公司正式季度指引。预告未经审计，不能机械年化；H1 EPS优先使用公司公告，未披露者由H1归母中值除以Q1股本推算。

\subsection{本轮最重要的估值含义}

\textbf{北京君正、德明利}补齐了存储观察池的预告缺口。北京君正H1归母中值约11.8亿元、扣非中值约11.5亿元，隐含Q2约8.6亿元，主营改善较清晰，但现价仍处高位，先做业绩验证；德明利H1归母中值约61.0亿元、扣非约60.5亿元，H1 EPS约27.13元，业绩已经兑现但价格先行，不能沿用预告前的旧分母直接追价。
\textbf{长电科技}是本轮先进封装中最干净的增量：归母与扣非中值分别约8.6和8.25亿元，非经常性占比约4.1\%，公司明确提到AI基础设施需求、订单和产能利用率。但一年位置和前期涨幅已经较高，结论是等待中报毛利率、资本开支回报和现金流，而不是把预告直接转化为新目标价。
\textbf{华天科技、通富微电}的归母利润不能直接当作主营分母。华天科技H1归母中值约8.0亿元、扣非约2.4亿元，非经常性占比约70\%；通富微电归母约17.0亿元、扣非约7.5亿元，非经常性占比约55.9\%。两者保留为产业趋势观察，正式估值必须等待中报解释公允价值、投资收益、毛利率和现金流。
\textbf{中航沈飞}的新增预告改变了原先“等待订单修复”的风险等级：H1归母约4.75亿元，同比预减约58.2\%，扣非约4.04亿元，同比预减约62.4\%。交付节奏和新型装备生产组织仍是约束，原先正向EPS分母不得继续作为升级依据。\textbf{航发动力}虽预增，但归母中值仅约1.48亿元、扣非约2.03亿元，绝对利润基数仍低，预告只支持订单交付验证，不支持估值重估。
\textbf{恺英网络、中国卫星}分别属于游戏经营改善和卫星型号验收进度改善。前者扣非中值约11.7亿元，仍有传奇IP和新品贡献，但非经常性占比约18.2\%；后者由亏转盈，扣非中值仅约0.29亿元，须观察H2合同履约和利润确认。两者均更新为“预告后验证”，不自动升级为正式核心。

\begin{riskbox}[增量数据边界]
本轮只补齐已被官方公告确认的9只标的；“未在新增清单”不代表公司没有预告。主报告的全市场母池仍以7月11日结构化普查为基线，后续若要把所有364家公司完全滚动到7月15日，需要重新抓取全量公告并重建行业/候选/估值矩阵。
\end{riskbox}
\sourcenote{data/earnings\_preview\_update\_20260715.json；sources/earnings-previews-20260715/。}
